\subsection{Proof of the Main Theorem}\label{subsec:proof}

Now we are ready to prove Theorem \ref{thm:main}.

\begin{proof}[Proof of Theorem \ref{thm:main}]
    We write
    \[C^b_n=\sum_{\mathcal{O}\in\mathcal{U}_n}\lvert \mathcal{O}\rvert\cdot r_b(\mathcal{O}; 0).\]
    From Lemma \ref{lem:orbit sizes}, $\lvert \mathcal{O}\rvert$ is always a power of two with exponent at least $s=s_2(n+1)-1$. We also know that $r_b(\mathcal{O}; 0)$ is always an integer and by definition, $r_b(\mathcal{O}; 0)$ is odd if and only if $\varepsilon^{\mathcal{O}}_0=1$. Then
    \[C^b_n\equiv \Big(\sum_{\mathcal{O}\in\mathcal{U}_n^{\text{min}}}\varepsilon^{\O}_0\Big)2^s\pmod{2^{s+1}},\]
    so it is equivalent to show that $\sum_{\mathcal{O}\in\mathcal{U}_n^{\text{min}}}\varepsilon^{\O}_0$ is odd.

    We use Theorem \ref{thm:orbit structure} to construct all minimal orbits of trees on $n$ vertices. By Corollary \ref{lem:reduction}, if we replace each complete binary tree of depth greater than $0$ by a single vertex, we get a new orbit of binary trees (on fewer vertices) which has the same sequence $\varepsilon^{\mathcal{O}}_m$, only reduced by a power of $\varepsilon_0$. This operation will be referred to as ``reduction.'' In Figure \ref{fig:oddcase}, have reduced the tree from Figure \ref{fig:minorbit} and obtained a new tree. Note that we have preserved the colors of the vertices, as we will need to refer to them later.
    
    \begin{figure}[h!]
    \centering
    \begin{tikzpicture}[
    level distance=1cm,
    level 1/.style={sibling distance=3cm},
    level 2/.style={sibling distance=2cm},
    level 3/.style={sibling distance=1.5cm},
    level 4/.style={sibling distance=0.75cm},
    level 5/.style={sibling distance=0.375cm},
    level 6/.style={sibling distance=0.1875cm}]
      \node[black] {}
        child {node[black] {}
          child {node[white] {}
            child {node[white] {}
            child {node[white] {}
            child {node[white] {}
            child {node[white] {}}
            child {node[white] {}}}
            child {node[white] {}
            child {node[white] {}}
            child {node[white] {}}}
            }
            child {node[white] {}
            child {node[white] {}
            child {node[white] {}}
            child {node[white] {}}}
            child {node[white] {}
            child {node[white] {}}
            child {node[white] {}}}
            }
            }
            child {node[white] {}
            child {node[white] {}
            child {node[white] {}
            child {node[white] {}}
            child {node[white] {}}}
            child {node[white] {}
            child {node[white] {}}
            child {node[white] {}}}
            }
            child {node[white] {}
            child {node[white] {}
            child {node[white] {}}
            child {node[white] {}}}
            child {node[white] {}
            child {node[white] {}}
            child {node[white] {}}}
            }
            }
            }
          child {node[white] {}
          child {node[white] {}}
          child {node[white] {}}}
        }
        child {node[black] {}
        child {node[black, xshift=1cm] {}
        child {node[white] {}}
        child {node[white] {}
            child {node[white] {}
            child {node[white] {}
            child {node[white] {}}
            child {node[white] {}}}
            child {node[white] {}
            child {node[white] {}}
            child {node[white] {}}}
            }
            child {node[white] {}
            child {node[white] {}
            child {node[white] {}}
            child {node[white] {}}}
            child {node[white] {}
            child {node[white] {}}
            child {node[white] {}}}
            }
            }}
        };
        
    \node[black, xshift=7cm] {}
        child {node[black] {}
          child {node[white] {}}
          child {node[white] {}
          }
        }
        child {node[black] {}
        child {node[black, xshift=1cm] {}
        child {node[white] {}}
        child {node[white] {}
            }
            }
        };
    \end{tikzpicture}
    \caption{The example tree from Figure \ref{fig:minorbit}, and its reduction on the right.}
    \label{fig:oddcase}
    \end{figure}
    
    We first consider the number of vertices remaining after reduction. If $n+1$ is odd, then there are $2s$ vertices remaining because there are $s$ nonempty complete binary trees being attached to $s$ vertices. If $n+1$ is even, then the number of remaining vertices is $2s+1$. In either case, let $k$ be the number of vertices of the trees of the minimal orbits after the reduction. With $\mathcal{U}_k$ as the set of orbits on $k$ vertices, let $f(\mathcal{O})$ denote the number of times an orbit $\mathcal{O}$ appears after reduction. Then we have
    \[\sum_{\mathcal{O}\in\mathcal{U}_n^{\text{min}}}\varepsilon^{\mathcal{O}}_0 = \varepsilon_0^{n-k}\sum_{\mathcal{O}\in \mathcal{U}_k}f(\mathcal{O})\varepsilon^{\mathcal{O}}_0.\]
    We want to determine the left hand side modulo $2$, so we can ignore all orbits for which $f(\mathcal{O})$ is even.
    
    In fact, we claim that those trees which remain (appear an odd number of times) are those representing minimal orbits of trees on $k$ vertices. Recall with Theorem \ref{thm:orbit structure} we assigned certain vertices as white and others as black, and we maintain this color scheme. After the reduction, we always have $s$ or $s+1$ white leaves attached to $s$ black vertices. See Table \ref{tab:reducedtrees} for an example, with all minimal orbits of trees on $14$ vertices shown with their reductions. In Table \ref{tab:reducedcount}, the counts of the resulting trees on $6$ vertices are recorded, and we can see that only three of them remain.
    
    \begin{table}[h]
    \centering
    \begin{tabular}{|c|c||c|c||c|c|} \hline
    Original tree & Reduction & Original Tree & Reduction & Original Tree & Reduction \\ \hline
    %1st
    \begin{tikzpicture}[
    level distance=1cm,
    level 1/.style={sibling distance=3cm},
    level 2/.style={sibling distance=1.5cm},
    level 3/.style={sibling distance=0.75cm},
    level 4/.style={sibling distance=0.375cm},
    scale = 0.5]
      \node[black] {}
        child {node[black] {}
          child {node[white, xshift=0.375cm] {}}
        }
        child {node[black] {}
        child {node[white] {}
        child {node[white] {}}
        child {node[white] {}}}
            child {node[white] {}
            child {node[white] {}
            child {node[white] {}}
            child {node[white] {}}}
            child {node[white] {}
            child {node[white] {}}
            child {node[white] {}}}}
        };
    \end{tikzpicture}
    &
    \begin{tikzpicture}[
    level distance=1cm,
    level 1/.style={sibling distance=3cm},
    level 2/.style={sibling distance=1.5cm},
    level 3/.style={sibling distance=0.75cm},
    level 4/.style={sibling distance=0.375cm},
    scale = 0.5]
        \node[xshift = 4cm, black] {}
        child {node[black] {}
          child {node[white, xshift=0.375cm] {}}
        }
        child {node[black] {}
        child {node[white] {}}
        child {node[white] {}}
        };
    \end{tikzpicture}
    
    &
    %7th
    \begin{tikzpicture}[
    level distance=1cm,
    level 1/.style={sibling distance=1.5cm},
    level 2/.style={sibling distance=1.5cm},
    level 3/.style={sibling distance=1.5cm},
    level 4/.style={sibling distance=0.75cm},
    level 5/.style={sibling distance=0.375cm},
    scale = 0.5]
      \node[black] {}
      child {node[white] {}}
        child {node[black] {}
        child {node[black, xshift=0.375cm] {}
        child {node[white] {}
        child {node[white] {}
        child {node[white] {}}
        child {node[white] {}}}
        child {node[white] {}
        child {node[white] {}}
        child {node[white] {}}}}
        child {node[white] {}
        child {node[white] {}}
        child {node[white] {}}}
        }
        };
    \end{tikzpicture}
    &
    \begin{tikzpicture}[
    level distance=1cm,
    level 1/.style={sibling distance=1.5cm},
    level 2/.style={sibling distance=1.5cm},
    level 3/.style={sibling distance=1.5cm},
    level 4/.style={sibling distance=0.75cm},
    level 5/.style={sibling distance=0.375cm},
    scale = 0.5]
        \node[black, xshift=4cm] {}
        child {node[white] {}}
        child {node[black] {}
        child {node[black, xshift = 0.375cm] {}
        child {node[white] {}}
        child {node[white] {}}
        }
        };
    \end{tikzpicture}
    
    &
    
    %13th
    \begin{tikzpicture}[
    level distance=1cm,
    level 1/.style={sibling distance=3.0cm},
    level 2/.style={sibling distance=1.5cm},
    level 3/.style={sibling distance=0.75cm},
    level 4/.style={sibling distance=0.375cm},
    level 5/.style={sibling distance=0.375cm},
    scale = 0.5]
      \node[black] {}
      child {node[white] {}
      child {node[white] {}}
      child {node[white] {}}
        }
        child {node[black] {}
        child {node[white] {}
        child {node[white] {}
        child {node[white] {}}
        child {node[white] {}}}
        child {node[white] {}
        child {node[white] {}}
        child {node[white] {}}}
        }
        child {node[black] {}
        child {node[white, xshift=-0.1875cm] {}}
        }
        };
    \end{tikzpicture}
    &
    \begin{tikzpicture}[
    level distance=1cm,
    level 1/.style={sibling distance=1.5cm},
    level 2/.style={sibling distance=1.5cm},
    level 3/.style={sibling distance=1.5cm},
    level 4/.style={sibling distance=0.75cm},
    level 5/.style={sibling distance=0.375cm},
    scale = 0.5]
        \node[black, xshift=4cm] {}
        child {node[white] {}}
        child {node[black] {}
        child {node[white] {}}
        child {node[black] {}
        child {node[white, xshift = -0.375cm] {}}
        }
        };
    \end{tikzpicture}
    \\ \hline
    
    %2nd
    \begin{tikzpicture}[
    level distance=1cm,
    level 1/.style={sibling distance=3cm},
    level 2/.style={sibling distance=1.5cm},
    level 3/.style={sibling distance=0.75cm},
    level 4/.style={sibling distance=0.375cm},
    scale = 0.5]
      \node[black] {}
        child {node[black] {}
          child {node[white, xshift=0.375cm] {}
          child {node[white] {}}
            child {node[white] {}}
            }
        }
        child {node[black] {}
        child {node[white] {}}
            child {node[white] {}
            child {node[white] {}
            child {node[white] {}}
            child {node[white] {}}}
            child {node[white] {}
            child {node[white] {}}
            child {node[white] {}}}}
        };
    \end{tikzpicture}
    &
    \begin{tikzpicture}[
    level distance=1cm,
    level 1/.style={sibling distance=3cm},
    level 2/.style={sibling distance=1.5cm},
    level 3/.style={sibling distance=0.75cm},
    level 4/.style={sibling distance=0.375cm},
    scale = 0.5]
        \node[xshift = 4cm, black] {}
        child {node[black] {}
          child {node[white, xshift=0.375cm] {}}
        }
        child {node[black] {}
        child {node[white] {}}
        child {node[white] {}}
        };
    \end{tikzpicture}
    
    &
    %8th
    \begin{tikzpicture}[
    level distance=1cm,
    level 1/.style={sibling distance=1.5cm},
    level 2/.style={sibling distance=1.5cm},
    level 3/.style={sibling distance=0.75cm},
    level 4/.style={sibling distance=0.75cm},
    level 5/.style={sibling distance=0.375cm},
    scale = 0.5]
      \node[black] {}
      child {node[white] {}
      child {node[white] {}
        child {node[white] {}}
        child {node[white] {}}}
        child {node[white] {}
        child {node[white] {}}
        child {node[white] {}}}
        }
        child {node[black] {}
        child {node[black, xshift=0.375cm] {}
        child {node[white] {}}
        child {node[white] {}
        child {node[white] {}}
        child {node[white] {}}}
        }
        };
    \end{tikzpicture}
    &
    \begin{tikzpicture}[
    level distance=1cm,
    level 1/.style={sibling distance=1.5cm},
    level 2/.style={sibling distance=1.5cm},
    level 3/.style={sibling distance=1.5cm},
    level 4/.style={sibling distance=0.75cm},
    level 5/.style={sibling distance=0.375cm},
    scale = 0.5]
        \node[black, xshift=4cm] {}
        child {node[white] {}}
        child {node[black] {}
        child {node[black, xshift = 0.375cm] {}
        child {node[white] {}}
        child {node[white] {}}
        }
        };
    \end{tikzpicture}
    
    &
    
    %14th
    \begin{tikzpicture}[
    level distance=1cm,
    level 1/.style={sibling distance=3.0cm},
    level 2/.style={sibling distance=1.5cm},
    level 3/.style={sibling distance=0.75cm},
    level 4/.style={sibling distance=0.375cm},
    level 5/.style={sibling distance=0.375cm},
    scale = 0.5]
      \node[black] {}
      child {node[white] {}
      child {node[white] {}
        child {node[white] {}}
        child {node[white] {}}}
        child {node[white] {}
        child {node[white] {}}
        child {node[white] {}}}
        }
        child {node[black] {}
        child {node[white] {}
        child {node[white] {}}
        child {node[white] {}}
        }
        child {node[black] {}
        child {node[white, xshift=-0.1875cm] {}}
        }
        };
    \end{tikzpicture}
    &
    \begin{tikzpicture}[
    level distance=1cm,
    level 1/.style={sibling distance=1.5cm},
    level 2/.style={sibling distance=1.5cm},
    level 3/.style={sibling distance=1.5cm},
    level 4/.style={sibling distance=0.75cm},
    level 5/.style={sibling distance=0.375cm},
    scale = 0.5]
        \node[black, xshift=4cm] {}
        child {node[white] {}}
        child {node[black] {}
        child {node[white] {}}
        child {node[black] {}
        child {node[white, xshift = -0.375cm] {}}
        }
        };
    \end{tikzpicture}
    \\ \hline
    
    %3rd
    \begin{tikzpicture}[
    level distance=1cm,
    level 1/.style={sibling distance=3cm},
    level 2/.style={sibling distance=1.5cm},
    level 3/.style={sibling distance=0.75cm},
    level 4/.style={sibling distance=0.375cm},
    scale = 0.5]
      \node[black] {}
        child {node[black] {}
          child {node[white, xshift=0.375cm] {}
          child {node[white] {}
            child {node[white] {}}
            child {node[white] {}}}
            child {node[white] {}
            child {node[white] {}}
            child {node[white] {}}}}
        }
        child {node[black] {}
        child {node[white] {}
        child {node[white] {}}
        child {node[white] {}}}
            child {node[white] {}}
        };
    \end{tikzpicture}
    &
    \begin{tikzpicture}[
    level distance=1cm,
    level 1/.style={sibling distance=3cm},
    level 2/.style={sibling distance=1.5cm},
    level 3/.style={sibling distance=0.75cm},
    level 4/.style={sibling distance=0.375cm},
    scale = 0.5]
        \node[xshift = 4cm, black] {}
        child {node[black] {}
          child {node[white, xshift=0.375cm] {}}
        }
        child {node[black] {}
        child {node[white] {}}
        child {node[white] {}}
        };
    \end{tikzpicture} 
    
    &
    %9th
    \begin{tikzpicture}[
    level distance=1cm,
    level 1/.style={sibling distance=1.5cm},
    level 2/.style={sibling distance=1.5cm},
    level 3/.style={sibling distance=1.5cm},
    level 4/.style={sibling distance=0.75cm},
    level 5/.style={sibling distance=0.375cm},
    scale = 0.5]
      \node[black] {}
      child {node[white] {}
      child {node[white] {}}
        child {node[white] {}}
        }
        child {node[black] {}
        child {node[black, xshift=0.375cm] {}
        child {node[white] {}
        child {node[white] {}
        child {node[white] {}}
        child {node[white] {}}}
        child {node[white] {}
        child {node[white] {}}
        child {node[white] {}}}}
        child {node[white] {}}
        }
        };
    \end{tikzpicture}
    &
    \begin{tikzpicture}[
    level distance=1cm,
    level 1/.style={sibling distance=1.5cm},
    level 2/.style={sibling distance=1.5cm},
    level 3/.style={sibling distance=1.5cm},
    level 4/.style={sibling distance=0.75cm},
    level 5/.style={sibling distance=0.375cm},
    scale = 0.5]
        \node[black, xshift=4cm] {}
        child {node[white] {}}
        child {node[black] {}
        child {node[black, xshift = 0.375cm] {}
        child {node[white] {}}
        child {node[white] {}}
        }
        };
    \end{tikzpicture}
    
    &
    
    %15th
    \begin{tikzpicture}[
    level distance=1cm,
    level 1/.style={sibling distance=3.0cm},
    level 2/.style={sibling distance=1.5cm},
    level 3/.style={sibling distance=0.75cm},
    level 4/.style={sibling distance=0.75cm},
    level 5/.style={sibling distance=0.75cm},
    scale = 0.5]
      \node[black] {}
      child {node[white] {}
      child {node[white] {}
        child {node[white] {}}
        child {node[white] {}}}
        child {node[white] {}
        child {node[white] {}}
        child {node[white] {}}}
        }
        child {node[black] {}
        child {node[white] {}
        }
        child {node[black] {}
        child {node[white, xshift=-0.1875cm] {}
        child {node[white] {}}
        child {node[white] {}}
        }
        }
        };
    \end{tikzpicture}
    &
    \begin{tikzpicture}[
    level distance=1cm,
    level 1/.style={sibling distance=1.5cm},
    level 2/.style={sibling distance=1.5cm},
    level 3/.style={sibling distance=1.5cm},
    level 4/.style={sibling distance=0.75cm},
    level 5/.style={sibling distance=0.375cm},
    scale = 0.5]
        \node[black, xshift=4cm] {}
        child {node[white] {}}
        child {node[black] {}
        child {node[white] {}}
        child {node[black] {}
        child {node[white, xshift = -0.375cm] {}}
        }
        };
    \end{tikzpicture}
    \\ \hline
    
    %4th
    \begin{tikzpicture}[
    level distance=1cm,
    level 1/.style={sibling distance=3cm},
    level 2/.style={sibling distance=3cm},
    level 3/.style={sibling distance=1.5cm},
    level 4/.style={sibling distance=0.75cm},
    level 5/.style={sibling distance=0.375cm},
    scale = 0.5]
      \node[black] {}
        child {node[black, xshift=0.75cm] {}
        child {node[white] {}}
        child {node[black] {}
        child {node[white] {}
        child {node[white] {}}
        child {node[white] {}}}
        child {node[white] {}
        child {node[white] {}
        child {node[white] {}}
        child {node[white] {}}}
        child {node[white] {}
        child {node[white] {}}
        child {node[white] {}}}}
        }
        };
    \end{tikzpicture}
    &
    \begin{tikzpicture}[
    level distance=1cm,
    level 1/.style={sibling distance=3cm},
    level 2/.style={sibling distance=3cm},
    level 3/.style={sibling distance=1.5cm},
    level 4/.style={sibling distance=0.75cm},
    level 5/.style={sibling distance=0.375cm},
    scale = 0.5]
        \node[black, xshift=4cm] {}
        child {node[black, xshift = 0.75cm] {}
        child {node[white] {}}
        child {node[black] {}
        child {node[white] {}}
        child {node[white] {}}
        }
        };
    \end{tikzpicture}
    
    &
    
    %10th
    \begin{tikzpicture}[
    level distance=1cm,
    level 1/.style={sibling distance=1.5cm},
    level 2/.style={sibling distance=1.5cm},
    level 3/.style={sibling distance=0.75cm},
    level 4/.style={sibling distance=0.75cm},
    level 5/.style={sibling distance=0.375cm},
    scale = 0.5]
      \node[black] {}
      child {node[white] {}
        }
        child {node[black] {}
        child {node[white] {}
        child {node[white] {}}
        child {node[white] {}}}
        child {node[black] {}
        child {node[white, xshift=-0.1875cm] {}
        child {node[white] {}
        child {node[white] {}}
        child {node[white] {}}}
        child {node[white] {}
        child {node[white] {}}
        child {node[white] {}}}}
        }
        };
    \end{tikzpicture}
    &
    \begin{tikzpicture}[
    level distance=1cm,
    level 1/.style={sibling distance=1.5cm},
    level 2/.style={sibling distance=1.5cm},
    level 3/.style={sibling distance=1.5cm},
    level 4/.style={sibling distance=0.75cm},
    level 5/.style={sibling distance=0.375cm},
    scale = 0.5]
        \node[black, xshift=4cm] {}
        child {node[white] {}}
        child {node[black] {}
        child {node[white] {}}
        child {node[black] {}
        child {node[white, xshift = -0.375cm] {}}
        }
        };
    \end{tikzpicture} & & \\ \hline
    
    %5th
    \begin{tikzpicture}[
    level distance=1cm,
    level 1/.style={sibling distance=3cm},
    level 2/.style={sibling distance=3cm},
    level 3/.style={sibling distance=1.5cm},
    level 4/.style={sibling distance=0.75cm},
    level 5/.style={sibling distance=0.375cm},
    scale = 0.5]
      \node[black] {}
        child {node[black, xshift=0.75cm] {}
        child {node[white] {}
        child {node[white] {}}
        child {node[white] {}}}
        child {node[black] {}
        child {node[white] {}}
        child {node[white] {}
        child {node[white] {}
        child {node[white] {}}
        child {node[white] {}}}
        child {node[white] {}
        child {node[white] {}}
        child {node[white] {}}}}
        }
        };
    \end{tikzpicture}
    &
    \begin{tikzpicture}[
    level distance=1cm,
    level 1/.style={sibling distance=3cm},
    level 2/.style={sibling distance=3cm},
    level 3/.style={sibling distance=1.5cm},
    level 4/.style={sibling distance=0.75cm},
    level 5/.style={sibling distance=0.375cm},
    scale = 0.5]
        \node[black, xshift=4cm] {}
        child {node[black, xshift = 0.75cm] {}
        child {node[white] {}}
        child {node[black] {}
        child {node[white] {}}
        child {node[white] {}}
        }
        };
    \end{tikzpicture}
    
    &
    
    %11th
    \begin{tikzpicture}[
    level distance=1cm,
    level 1/.style={sibling distance=1.5cm},
    level 2/.style={sibling distance=1.5cm},
    level 3/.style={sibling distance=0.75cm},
    level 4/.style={sibling distance=0.375cm},
    level 5/.style={sibling distance=0.375cm},
    scale = 0.5]
      \node[black] {}
      child {node[white] {}
        }
        child {node[black] {}
        child {node[white] {}
        child {node[white] {}
        child {node[white] {}}
        child {node[white] {}}}
        child {node[white] {}
        child {node[white] {}}
        child {node[white] {}}}
        }
        child {node[black] {}
        child {node[white, xshift=-0.1875cm] {}
        child {node[white] {}}
        child {node[white] {}}
        }
        }
        };
    \end{tikzpicture}
    &
    \begin{tikzpicture}[
    level distance=1cm,
    level 1/.style={sibling distance=1.5cm},
    level 2/.style={sibling distance=1.5cm},
    level 3/.style={sibling distance=1.5cm},
    level 4/.style={sibling distance=0.75cm},
    level 5/.style={sibling distance=0.375cm},
    scale = 0.5]
        \node[black, xshift=4cm] {}
        child {node[white] {}}
        child {node[black] {}
        child {node[white] {}}
        child {node[black] {}
        child {node[white, xshift = -0.375cm] {}}
        }
        };
    \end{tikzpicture} & & \\ \hline
    
    %6th
    \begin{tikzpicture}[
    level distance=1cm,
    level 1/.style={sibling distance=3cm},
    level 2/.style={sibling distance=3cm},
    level 3/.style={sibling distance=1.5cm},
    level 4/.style={sibling distance=0.75cm},
    level 5/.style={sibling distance=0.375cm},
    scale = 0.5]
      \node[black] {}
        child {node[black, xshift=0.75cm] {}
        child {node[white] {}
        child {node[white] {}
        child {node[white] {}}
        child {node[white] {}}}
        child {node[white] {}
        child {node[white] {}}
        child {node[white] {}}}}
        child {node[black] {}
        child {node[white] {}}
        child {node[white] {}
        child {node[white] {}}
        child {node[white] {}}}
        }
        };
    \end{tikzpicture}
    &
    \begin{tikzpicture}[
    level distance=1cm,
    level 1/.style={sibling distance=3cm},
    level 2/.style={sibling distance=3cm},
    level 3/.style={sibling distance=1.5cm},
    level 4/.style={sibling distance=0.75cm},
    level 5/.style={sibling distance=0.375cm},
    scale = 0.5]
        \node[black, xshift=4cm] {}
        child {node[black, xshift = 0.75cm] {}
        child {node[white] {}}
        child {node[black] {}
        child {node[white] {}}
        child {node[white] {}}
        }
        };
    \end{tikzpicture}
    
    &
    
    %12th
    \begin{tikzpicture}[
    level distance=1cm,
    level 1/.style={sibling distance=1.5cm},
    level 2/.style={sibling distance=0.75cm},
    level 3/.style={sibling distance=0.75cm},
    level 4/.style={sibling distance=0.75cm},
    level 5/.style={sibling distance=0.375cm},
    scale = 0.5]
      \node[black] {}
      child {node[white] {}
      child {node[white] {}}
      child {node[white] {}}
        }
        child {node[black] {}
        child {node[white] {}
        }
        child {node[black] {}
        child {node[white, xshift=-0.1875cm] {}
        child {node[white] {}
        child {node[white] {}}
        child {node[white] {}}}
        child {node[white] {}
        child {node[white] {}}
        child {node[white] {}}}}
        }
        };
    \end{tikzpicture}
    &
    \begin{tikzpicture}[
    level distance=1cm,
    level 1/.style={sibling distance=1.5cm},
    level 2/.style={sibling distance=1.5cm},
    level 3/.style={sibling distance=1.5cm},
    level 4/.style={sibling distance=0.75cm},
    level 5/.style={sibling distance=0.375cm},
    scale = 0.5]
        \node[black, xshift=4cm] {}
        child {node[white] {}}
        child {node[black] {}
        child {node[white] {}}
        child {node[black] {}
        child {node[white, xshift = -0.375cm] {}}
        }
        };
    \end{tikzpicture} & & \\ \hline
    \end{tabular}
    \caption{Reduction of minimal orbits of trees on $14$ vertices to orbits of trees on $6$ vertices.}
    \label{tab:reducedtrees}
\end{table}

\begin{table}[h]
    \centering
    \begin{tabular}{|c|c|}\hline
    Reduced Tree & Count (=$f(\O)$) \\ \hline
    \begin{tikzpicture}[
    level distance=1cm,
    level 1/.style={sibling distance=3cm},
    level 2/.style={sibling distance=1.5cm},
    level 3/.style={sibling distance=0.75cm},
    level 4/.style={sibling distance=0.375cm},
    scale = 0.5]
        \node[xshift = 4cm, black] {}
        child {node[black] {}
          child {node[white, xshift=0.375cm] {}}
        }
        child {node[black] {}
        child {node[white] {}}
        child {node[white] {}}
        };
    \end{tikzpicture} & $3=\dfrac{3!\cdot 4}{2^3}$ \\ \hline
    
    \begin{tikzpicture}[
    level distance=1cm,
    level 1/.style={sibling distance=3cm},
    level 2/.style={sibling distance=3cm},
    level 3/.style={sibling distance=1.5cm},
    level 4/.style={sibling distance=0.75cm},
    level 5/.style={sibling distance=0.375cm},
    scale = 0.5]
        \node[black, xshift=4cm] {}
        child {node[black, xshift = 0.75cm] {}
        child {node[white] {}}
        child {node[black] {}
        child {node[white] {}}
        child {node[white] {}}
        }
        };
    \end{tikzpicture} & $3=\dfrac{3!\cdot 4}{2^3}$ \\ \hline
    
    \begin{tikzpicture}[
    level distance=1cm,
    level 1/.style={sibling distance=1.5cm},
    level 2/.style={sibling distance=1.5cm},
    level 3/.style={sibling distance=1.5cm},
    level 4/.style={sibling distance=0.75cm},
    level 5/.style={sibling distance=0.375cm},
    scale = 0.5]
        \node[black, xshift=4cm] {}
        child {node[white] {}}
        child {node[black] {}
        child {node[black, xshift = 0.375cm] {}
        child {node[white] {}}
        child {node[white] {}}
        }
        };
    \end{tikzpicture} & $3=\dfrac{3!\cdot 4}{2^3}$ \\ \hline
    
    \begin{tikzpicture}[
    level distance=1cm,
    level 1/.style={sibling distance=1.5cm},
    level 2/.style={sibling distance=1.5cm},
    level 3/.style={sibling distance=1.5cm},
    level 4/.style={sibling distance=0.75cm},
    level 5/.style={sibling distance=0.375cm},
    scale = 0.5]
        \node[black, xshift=4cm] {}
        child {node[white] {}}
        child {node[black] {}
        child {node[white] {}}
        child {node[black] {}
        child {node[white, xshift = -0.375cm] {}}
        }
        };
    \end{tikzpicture} & $6=\dfrac{3!\cdot 8}{2^3}$ \\ \hline
    \end{tabular}
    \caption{Counts of the resulting reduced trees.}
    \label{tab:reducedcount}
\end{table}
    
    
    Let $M$ be the size of the orbit of a tree on $k$ vertices obtained by reduction. In the case where $k=2s$, we claim that the number of times a given orbit $\O$ appears is given by $\frac{s!M}{2^s}$. We can prove this by counting the number of trees on $n$ vertices to which this corresponds: we have $s!$ ways of placing complete binary trees on the $s$ white vertices and $M$ different ways to determine the underlying structure, which will count every corresponding tree on $n$ vertices $2^s$ times because each such tree has $2^s$ symmetries.
    
    Now we determine when this quantity is odd. We have
    \[\xi_2\left(\frac{s!M}{2^s}\right)=\xi_2(s!)+\xi_2(M)-\xi_2(2^s)=s-s_2(s)+\xi_2(M)-s=\xi_2(M)-s_2(s).\]
    This is odd if and only if $\xi_2(M)=s_2(s)$. Since $s_2(s)=s_2(2s+1)-1=s_2(k)-1$, this occurs if and only if we have a minimal orbit.
    
    Similarly, when $k=2s+1$, the number of times that a tree appears is given by $\frac{(s+1)!M}{2^s}$, which is odd only when we have a minimal orbit for the same reason.
    
    Up to now, we have shown that if $f(\mathcal{O})\geq 1$, i.e. an orbit of trees on $k$ vertices actually appears after the reduction operation, then $f(\mathcal{O})$ is even if $\mathcal{O}\not \in\mathcal{U}_k^{\text{min}}$. If $\mathcal{O}\in \mathcal{U}_k^{\text{min}}$, then $\mathcal{O}$ will be counted an odd number of times if we know how to assign the nodes as white/black to ``expand'' the tree back to the original size, though this still leaves out the possibility of multiple or no white/black assignments for $\O$. To finish the claim, we show that if $\mathcal{O}\in\mathcal{U}_k^{\text{min}}$, then there is a unique assignment of black and white nodes to $\O$ so that some orbits in $\mathcal{U}^{\text{min}}_n$ reduce to $\O$ with that color assignment. We know that nodes with no children after reduction must be white and those with at least one child must be black after reduction, which shows uniqueness. Using Theorem \ref{thm:orbit structure}, we may directly verify that trees in an orbit of $\mathcal{U}_k^{\text{min}}$ will always have a number of white nodes which is equal to or one more than the number of black nodes, so this coloring is always valid.
    
    With the claim, we now have
    \[\sum_{\mathcal{O}\in\mathcal{U}_n^{\text{min}}}\varepsilon^{\mathcal{O}}_0= \varepsilon_0^{n-k}\sum_{\mathcal{O}\in \mathcal{U}_k}f(\mathcal{O})\varepsilon^{\mathcal{O}}_0 \equiv\varepsilon^{n-k}_0\sum_{\mathcal{O}\in\mathcal{U}_k^{\text{min}}}\varepsilon^{\mathcal{O}}_0\pmod{2}\]
    so we may reduce to the case of $k$ vertices. We always have $k < n$ if $n\geq 2$ because $n+1\geq 3$, meaning that minimal orbits contain complete binary trees of depth greater than $1$ and the reduction operation always removes a positive number of vertices.
    
    For our base cases, we will eventually reduce to $n=1$ or $n=2$, where the statement is directly verified. Specifically, we have
    \begin{align*}
        \sum_{\mathcal{O}\in\mathcal{U}_1^{\text{min}}}\varepsilon^{\mathcal{O}}_0 &= \varepsilon_0, \\
        \sum_{\mathcal{O}\in\mathcal{U}_2^{\text{min}}}\varepsilon^{\mathcal{O}}_0 &= \varepsilon_0(\varepsilon_0+\varepsilon_1).
    \end{align*}
    The assumptions on $b$ are precisely that $\varepsilon_0=1$ and $\varepsilon_1=0$, so this is odd, as desired.
    
    Note that the reduction to a smaller case preserves parity, so we reduce to the $n=1$ or $n=2$ case according to the parity of $n$ at the start. Thus we actually get a stronger result: if we relax the condition $4\mid \Delta b$ to $2\mid \Delta b$, then $\xi_2(C^b_n)=\xi_2(C_n)$ holds for odd $n$, while $\xi_2(C^b_n) > \xi_2(C_n)$ for even $n$ if $4\nmid \Delta b$.
\end{proof}

An \cite{an2010combinatorial} previously determined necessary and sufficient conditions, conjectured by Konvalinka \cite{konvalinka2007divisibility}, under which a polynomial weight function $b(x)=b_0+b_1x+b_2x^2+\dotsb\in \mathbb{Z}[x]$ satisfies $\xi_2(C^b_n)=\xi_2(C_n)$. Inspired by his results, we make the following conjecture.
\begin{conjecture}
\label{conj:maingeneralization}
If $b: \Z_{\geq 0}\to\Z$ satisfies 
\begin{enumerate}
    \item $b(0)$ is odd,
    \item $2^{n-s_2(n)}\mid (\Delta^nb)(x)$ for all $n\geq 2$ and $x\in \mathbb{Z}_{\geq 0}$,
    \item $b(0)\equiv b(1)\pmod{4}$,
\end{enumerate}
then $\xi_2(C_n^b)=\xi_2(C_n)=s_2(n+1)-1.$
\end{conjecture}
We would be particularly interested in seeing a combinatorial proof of Conjecture \ref{conj:maingeneralization} as An's overall strategy was very computational in nature.